Integral homology at coefficient two supplies the remaining family.
Its character sees homology modulo eight, where the square of a lift
of multiplier minus one records the necessary cyclic correction.

\begin{theorem}\label{thm:main-complete-quadratic}
For every odd prime \(p\), there is a continuous invariant
character \(\xi_{2,p}:M\to\F_2\), determined by compact homology
modulo eight of the degree-\(p\) cyclic cover with monodromies
\(1,1,-2\), such that \(\nu_p=\xi_{2,p}|_N\in X_2\) and
\[
 \operatorname{tr}(\nu_p)
     =\operatorname{Inf}_U^A
          \bigl(ac_p+\varepsilon_p(\delta_2+\delta_p)\bigr),
 \qquad \nu_p|_K=\kappa_{2,p}.
\]
The character is defined for every positive peripheral action with
the hypotheses of Theorem~\ref{thm:main-odd-pair-lifts}.
Consequently the original quadratic restriction is surjective, with
an explicit linear section:
\[
 0\longrightarrow\operatorname{Hom}_{\rm cts}(B,\F_2)
 \longrightarrow X_2\xrightarrow{\ \mathrm{res}\ }Y_2
 \longrightarrow0.
\]
One section sends the basis elements \(\eta_3,\psi_p,\kappa_{2,p},
\kappa_{\ell,q}\) to \(\eta_3,\psi_p,\nu_p,\xi_{\ell,q}\),
respectively, with \(\ell<q\). The companion character of the same
integral homology construction satisfies
\[
 \widetilde\psi_p=\rho_p^I/2\pmod2
                 \quad\text{on }M\text{ for every odd }p.
\]
At \(p=3\), the lift \(\nu_3\)
equals the cubic CM lift \(\tau+\psi_3\)
defined below; use this latter section \(s\) in the simultaneous
normalization, and put \(\mathcal S=s(Y_2)\). Then
\[
               X_2=\operatorname{Hom}_{\rm cts}(B,\F_2)
                                      \oplus\mathcal S.
\]
The two inflation kernels on \(H^2(A,\F_2)\), towards
\(H^2(G,\F_2)\) and towards \(\operatorname{Br}(\Q)[2]\),
are equal.
\end{theorem}

Every arithmetic invariant quadratic character therefore extends.
The unresolved part of the original character problem is exactly
\(\operatorname{Hom}_{\rm cts}(B,\F_2)\): surjectivity does not
prove that this kernel vanishes. The complete proofs are in
Sections~\ref{sec:coefficient-two-cross-cups} and
\ref{sec:companion-kummer-equality} and
\ref{sec:complete-quadratic-restriction}.
